\documentclass[a4paper,12pt]{article}
\usepackage[T2A]{fontenc}
\usepackage[utf8]{inputenc}
\usepackage[russian]{babel}
\usepackage{amsmath,amssymb,mathtools}
\usepackage{icomma}
\usepackage[margin=25mm,top=23mm,bottom=24mm]{geometry}
\usepackage{microtype}
\usepackage{tikz}
\usetikzlibrary{arrows.meta,positioning,shapes.geometric,calc}
\usepackage{tabularx,booktabs,array}
\usepackage{enumitem}
\usepackage{xcolor}
\usepackage{fancyhdr}
\usepackage{setspace}
\usepackage{hyperref}
\usepackage{titlesec}
\definecolor{accent}{HTML}{087F86}
\definecolor{darkaccent}{HTML}{08636A}
\definecolor{textgray}{HTML}{34393D}
\definecolor{muted}{HTML}{677176}
\hypersetup{colorlinks=true,linkcolor=darkaccent,urlcolor=darkaccent}
\setlength{\parindent}{0pt}
\setlength{\parskip}{5pt plus 1pt minus 1pt}
\setstretch{1.16}
\setlist[enumerate]{leftmargin=*,itemsep=11pt,topsep=5pt}
\setlist[itemize]{leftmargin=*,itemsep=4pt,topsep=4pt}
\titleformat{\section}{\Large\color{darkaccent}\normalfont}{\thesection.}{0.65em}{}
\titleformat{\subsection}{\normalsize\color{darkaccent}\normalfont}{}{0em}{}
\titlespacing*{\section}{0pt}{19pt}{10pt}
\titlespacing*{\subsection}{0pt}{14pt}{7pt}
\pagestyle{fancy}
\fancyhf{}
\fancyhead[L]{\footnotesize\color{muted}Аннуитетные кредиты}
\fancyhead[R]{\footnotesize\color{muted}Анастасия Павлова}
\fancyfoot[C]{\footnotesize\color{muted}\thepage}
\renewcommand{\headrulewidth}{0.35pt}
\renewcommand{\headrule}{\hbox to\headwidth{\color{accent}\leaders\hrule height \headrulewidth\hfill}}
\setlength{\headheight}{23pt}
\newcommand{\focus}[1]{\textcolor{darkaccent}{#1}}
\begin{document}
\thispagestyle{empty}
\begin{center}
\vspace*{23mm}
{\large\color{darkaccent} МАТЕМАТИКА}\par
\vspace{13mm}
{\LARGE Чек-лист}\par
\vspace{9mm}
{\LARGE\color{accent} Аннуитетная схема\\[3pt] погашения кредита}\par
\vspace{14mm}
{\normalsize Подготовка к ЕГЭ по профильной математике}\par
\vfill
{\large Анастасия Павлова}\par
\vspace{12mm}
{\normalsize 2 ноября 2026 г.}\par
\vspace{17mm}
{\small Лёвин Артём Александрович\\эксклюзивно для Анастасии Павловой}\par
\vspace{15mm}
\begin{tikzpicture}
\draw[accent,line width=0.65pt] (0,0) .. controls (1.3,0.6) and (2.4,-0.5) .. (3.9,0)
    .. controls (5.2,0.55) and (6.25,-0.4) .. (7.6,0);
\end{tikzpicture}
\vspace*{6mm}
\end{center}
\clearpage

\section{Как распознать аннуитетную схему}
\focus{Главный признак:} кредит погашается \textbf{равными платежами} через равные промежутки времени. При ежегодном начислении процентов и ежегодных платежах одна и та же сумма \(x\) выплачивается в конце каждого года.

На каждом этапе порядок одинаковый:
\[
\boxed{\text{Остаток долга}}
\ \xrightarrow{\text{начислить проценты}}\
\boxed{\text{Долг с процентами}}
\ \xrightarrow{\text{вычесть }x}\
\boxed{\text{Новый остаток}}.
\]
По окончании последнего года \textbf{остаток долга равен нулю}.

\section{Обозначения и повышающий коэффициент}
\begin{tabularx}{\linewidth}{@{}p{27mm}X@{}}
\toprule
\(S\) & сумма кредита в начале срока; \\
\(P\) & годовая процентная ставка (например, \(16\)); \\
\(p=P/100\) & та же ставка в виде десятичной дроби; \\
\(k=1+p\) & повышающий коэффициент; \\
\(n\) & число лет (и равных ежегодных платежей); \\
\(x\) & размер каждого ежегодного платежа. \\
\bottomrule
\end{tabularx}

Если ставка \(16\%\), то \(p=0{,}16\), \(k=1{,}16\). Долг \(D\) после начисления процентов превращается в \(Dk\). После платежа \(x\) остаётся
\[
\boxed{D_{\text{новый}}=Dk-x}.
\]
\textbf{Важно:} проценты каждого следующего года начисляют на оставшийся долг, а не на первоначальную сумму \(S\).

\section{Таблица для трёх лет}
Введём \(D_0=S\). После каждого года будем записывать остаток как \(D_1,D_2,D_3\).

\begingroup\small\renewcommand{\arraystretch}{1.45}
\begin{tabularx}{\linewidth}{@{}c>{\centering\arraybackslash}p{20mm}>{\centering\arraybackslash}p{22mm}>{\centering\arraybackslash}p{24mm}cX@{}}
\toprule
Год & Начало года & Проценты & Долг до платежа & Платёж & Остаток \\
\midrule
1 & \(S\) & \(Sp\) & \(Sk\) & \(x\) & \(D_1=Sk-x\) \\
2 & \(D_1\) & \(D_1p\) & \(D_1k\) & \(x\) & \(D_2=D_1k-x\) \\
3 & \(D_2\) & \(D_2p\) & \(D_2k\) & \(x\) & \(D_3=D_2k-x=0\) \\
\bottomrule
\end{tabularx}
\endgroup

\clearpage
\section{Как вывести формулу платежа}
Последовательно раскроем записи остатков для трёхлетнего кредита:
\begin{align*}
D_1&=Sk-x,\\
D_2&=(Sk-x)k-x=Sk^2-xk-x,\\
D_3&=(Sk^2-xk-x)k-x\\
   &=Sk^3-xk^2-xk-x.
\end{align*}
Кредит полностью погашен, поэтому \(D_3=0\):
\begin{align*}
Sk^3-x(k^2+k+1)&=0,\\
x(k^2+k+1)&=Sk^3,\\
\boxed{x=\frac{Sk^3}{k^2+k+1}}.&
\end{align*}
При вынесении \(x\) за скобки появляются степени \(k^2,k^1,k^0\). Последняя степень \(k^0=1\), поэтому последнее слагаемое в знаменателе --- единица.

\section{Формула на любое число лет}
Для \(n\) ежегодных равных платежей:
\[
\boxed{Sk^n-x\bigl(k^{n-1}+k^{n-2}+\ldots+k+1\bigr)=0},
\]
\[
\boxed{x=\frac{Sk^n}{k^{n-1}+k^{n-2}+\ldots+k+1}}.
\]
\focus{Закономерность:} старшая степень в числителе равна числу лет, а в знаменателе степени убывают от \(n-1\) до нуля.

\begin{tabularx}{\linewidth}{@{}cX@{}}
\toprule
\(n\) & Ежегодный платёж \(x\) \\
\midrule
2 & \(\displaystyle\frac{Sk^2}{k+1}\) \\[5pt]
3 & \(\displaystyle\frac{Sk^3}{k^2+k+1}\) \\[5pt]
4 & \(\displaystyle\frac{Sk^4}{k^3+k^2+k+1}\) \\[5pt]
5 & \(\displaystyle\frac{Sk^5}{k^4+k^3+k^2+k+1}\) \\
\bottomrule
\end{tabularx}

Формулы выводятся из таблицы или из повторяющегося правила \(D_j=kD_{j-1}-x\). Если \(n\) невелико, удобно сразу записать требуемую сумму степеней.

\clearpage
\section{Общая сумма, переплата и её процент}
Пусть все \(n\) платежей равны \(x\). Тогда
\begin{align*}
\text{Общая сумма платежей (ОСП)}&=nx,\\
\text{Переплата}&=nx-S,\\
\text{Процент переплаты}&=\frac{nx-S}{S}\cdot100\%.
\end{align*}
Это три \emph{разные} величины. Процентная ставка \(P\) обозначает начисление за один год, а процент переплаты показывает итоговую переплату за весь срок относительно исходного кредита.

\subsection{Пример 1. Кредит на два года}
Кредит \(810\) тыс. руб. выдан под \(16\%\) годовых и погашается двумя одинаковыми платежами. Найдите общую сумму платежей в рублях.

\textbf{Шаг 1.} \(S=810\) (тыс. руб.), \(n=2\), \(k=1{,}16\).

\textbf{Шаг 2.} Находим один платёж:
\[
 x=\frac{810\cdot 1{,}16^2}{1{,}16+1}
  =\frac{810\cdot 1{,}3456}{2{,}16}
  =504{,}6\quad\text{(тыс. руб.)}.
\]
\textbf{Шаг 3.} Находим сумму двух платежей:
\[
\text{ОСП}=2\cdot504{,}6=1009{,}2\quad\text{(тыс. руб.)}.
\]
\textbf{Шаг 4.} Переводим в рубли:
\[
\boxed{\text{ОСП}=1\,009\,200\ \text{руб.}}
\]
\textbf{Проверка:} остаток после первого года равен \(810\cdot1{,}16-504{,}6=435\) тыс. руб. После второго: \(435\cdot1{,}16-504{,}6=0\).

\subsection{Пример 2. Проверка таблицей}
Кредит \(91\) тыс. руб. на три года под \(20\%\) годовых.
\[
 x=\frac{91\cdot1{,}2^3}{1{,}2^2+1{,}2+1}=43{,}2\quad\text{(тыс. руб.)}.
\]
Остатки после платежей:
\[
91\cdot1{,}2-43{,}2=66,\qquad
66\cdot1{,}2-43{,}2=36,\qquad
36\cdot1{,}2-43{,}2=0.
\]

\clearpage
\section{Типичные ошибки и самопроверка}
\begin{enumerate}[itemsep=7pt]
\item \textbf{Перепутали схемы.} Слова «равными платежами» указывают на аннуитет. В схеме с погашением долга равными частями платежи обычно различаются.
\item \textbf{Начисляете проценты на старую сумму.} Каждый раз используется остаток после предыдущего платежа.
\item \textbf{Потеряли слагаемое \(1\).} При \(n=3\) в знаменателе именно \(k^2+k+1\), при \(n=5\) --- пять слагаемых от \(k^4\) до \(1\).
\item \textbf{Приняли \(x\) за ОСП.} \(x\) --- один платёж; общая сумма \(nx\).
\item \textbf{Забыли единицы.} Если \(S\) задано в тысячах рублей, платёж и переплата сначала получаются также в тысячах рублей.
\item \textbf{Не проверили погашение.} После последнего платежа долг обязан стать равным нулю.
\end{enumerate}

\subsection{Маршрут решения задачи}
\begin{enumerate}[itemsep=4pt]
\item Найдите в условии признак одинаковых платежей.
\item Запишите \(S\), \(P\), \(n\) и вычислите \(k=1+P/100\).
\item Составьте уравнение \(Sk^n=x(k^{n-1}+\ldots+1)\) и найдите \(x\).
\item Вычислите именно то, о чём спрашивают: \(x\), ОСП, переплату или её процент.
\item Проверьте единицы измерения и при необходимости остаток последнего года.
\end{enumerate}

\section{Границы применимости}
Эти формулы предполагают \textbf{одинаковую годовую ставку}, начисление процентов перед каждым ежегодным платежом и \textbf{равные} ежегодные платежи. При меняющейся ставке или особом графике платежей нужно отдельно составлять расчёт по годам.

\clearpage
\thispagestyle{plain}
\begin{center}
{\Large\color{darkaccent} Тренировочный блок}\par
\vspace{3pt}
{\small Путь А \(\cdot\) 10 заданий \(\cdot\) от простого к сложному}
\end{center}
\vspace{4mm}
Решайте последовательно. При необходимости указывайте промежуточные остатки долга и единицы измерения.

\begin{enumerate}[itemsep=12pt]
\item Найдите повышающий коэффициент \(k\) для годовых ставок \(8\%\), \(16\%\) и \(20\%\).
\item В начале года долг составляет \(100\) тыс. руб. За год он увеличивается на \(10\%\), затем вносится платёж \(40\) тыс. руб. Найдите остаток долга после платежа.
\item Кредит \(210\) тыс. руб. под \(10\%\) годовых погашается двумя равными ежегодными платежами. Найдите размер платежа и остаток после первого года.
\item Кредит \(200\) тыс. руб. под \(20\%\) годовых полностью погашают тремя равными ежегодными платежами \(x\) тыс. руб. Составьте уравнение, выражающее нулевой остаток после третьего года. Решать его не требуется.
\item Для кредита \(S\) сроком на три года с коэффициентом \(k\) выведите формулу ежегодного равного платежа \(x\).
\item Запишите формулу ежегодного равного платежа \(x\) для кредита \(S\), выдаваемого на пять лет с постоянным повышающим коэффициентом \(k\).
\item Кредит \(182\) тыс. руб. взят на три года под \(20\%\) годовых. Найдите ежегодный равный платёж.
\item Кредит \(110\) тыс. руб. под \(20\%\) годовых погашают двумя равными ежегодными платежами. Найдите общую сумму платежей и переплату \textbf{в рублях}.
\item Кредит погашен двумя ежегодными платежами по \(75\) тыс. руб. Ставка составляет \(25\%\) годовых. Найдите первоначальную сумму кредита.
\item Кредит \(464{,}1\) тыс. руб. выдан на четыре года под \(10\%\) годовых. Найдите ежегодный платёж, общую сумму платежей, переплату и процент переплаты. Последний ответ округлите до десятых процента.
\end{enumerate}

\clearpage
\thispagestyle{plain}
\begin{center}
{\Large\color{darkaccent} Ответы}\par
\vspace{3pt}
{\small К тренировочному блоку}
\end{center}
\vspace{6mm}
\renewcommand{\arraystretch}{1.55}
\begin{tabularx}{\linewidth}{@{}cX@{}}
\toprule
№ & Ответ \\
\midrule
1 & \(1{,}08;\ 1{,}16;\ 1{,}2\). \\
2 & \(70\) тыс. руб. \\
3 & Платёж \(121\) тыс. руб.; остаток \(110\) тыс. руб. \\
4 & \(\bigl((240-x)\cdot1{,}2-x\bigr)\cdot1{,}2-x=0\). \\
5 & \(\displaystyle x=\frac{Sk^3}{k^2+k+1}\). \\[7pt]
6 & \(\displaystyle x=\frac{Sk^5}{k^4+k^3+k^2+k+1}\). \\[7pt]
7 & \(86{,}4\) тыс. руб. \\
8 & ОСП \(144\,000\) руб.; переплата \(34\,000\) руб. \\
9 & \(108\) тыс. руб. \\
10 & \(x=146{,}41\) тыс. руб.; ОСП \(585{,}64\) тыс. руб.; переплата \(121{,}54\) тыс. руб.; \(26{,}2\%\). \\
\bottomrule
\end{tabularx}

\vspace{12mm}
\section*{Итоговая памятка}
\begin{tabularx}{\linewidth}{@{}p{54mm}X@{}}
\toprule
Признак или вопрос & Что делать \\
\midrule
Равные ежегодные платежи & Применяйте аннуитетную схему. \\
Проценты за год & Умножайте текущий долг на \(k=1+P/100\). \\
Остаток после года & Вычитайте \(x\) из долга с процентами. \\
Найти один платёж & Используйте формулу с \(Sk^n\). \\
Найти всю выплату & Вычислите \(nx\). \\
Найти переплату & Вычислите \(nx-S\). \\
Найти процент переплаты & Вычислите \(\frac{nx-S}{S}\cdot100\%\). \\
\bottomrule
\end{tabularx}
\end{document}
